\chapter*{Đối tượng và phương pháp nghiên cứu}\addcontentsline{toc}{chapter}{Đối tượng và phương pháp nghiên cứu}

\textbf{Đối tượng và phạm vi nghiên cứu}

Kết quả xổ số kiến thiết Miền Bắc, bao gồm toàn bộ các số trúng thưởng được công bố trong mỗi ngày quay số. 
Giải thưởng và thể lệ trúng thưởng được thu thập đúng theo công bố của công ty xổ số kiến thiết.

Dữ liệu được thu thập của xổ số kiến thiết Miền Bắc, do Miền Bắc sẽ chỉ có 1 kết quả duy nhất trong ngày, đảm bảo về tính toàn vẹn của thời gian (Miền Nam sẽ có nhiều kết quả trong ngày). Thời gian của kết quả đầy đủ từ năm 2007 đến hiện tại. 
Trong đó, giai đoạn hết năm 2022 được sử dụng làm thông tin lịch sử cho quá 
trình xây dựng mô hình. Giai đoạn 2023--2024 được sử dụng cho đánh giá và lựa 
chọn chiến lược, trong khi giai đoạn 2025--2026 được giữ lại làm tập kiểm định 
cuối cùng ngoài mẫu.

Mỗi ngày quay số bao gồm toàn bộ 27 kết quả thuộc các hạng giải từ 
giải đặc biệt đến giải bảy. Các kết quả được chuẩn hóa để phục vụ phân tích 
theo từng vị trí chữ số và theo phần hậu tố của số trúng thưởng.
	
Tần suất quan sát là theo ngày. Thứ tự thời gian của dữ liệu được giữ 
nguyên trong toàn bộ quá trình xây dựng và đánh giá mô hình; không thực hiện 
chia ngẫu nhiên dữ liệu huấn luyện và kiểm định.
	

\textbf{Phương pháp nghiên cứu}

\begin{itemize}
	\item \textbf{Kiểm định tính ngẫu nhiên:} nghiên cứu phân tích phân phối chữ 
	số, tính đồng đều, tính độc lập giữa các vị trí, sự phụ thuộc Markov, phụ 
	thuộc theo độ trễ và mức độ ổn định của phân phối theo thời gian. Khi cần 
	thiết, kích thước hiệu ứng và hiệu chỉnh cho nhiều phép kiểm định được xem xét 
	song song với giá trị $p$.
	
	\item \textbf{Mô hình dự báo:} các phương pháp thống kê, mô hình xác 
	suất và mô hình học máy được sử dụng để kiểm định tính ngẫu nhiên và đánh giá 
	khả năng dự báo kết quả xổ số. Các phương pháp chính bao gồm kiểm định phân 
	phối và tính độc lập, mô hình tần suất, Rolling Frequency, mô hình Markov, 
	Random Forest, XGBoost và CatBoost.
	
	\item \textbf{Đánh giá mô hình:} hiệu quả dự báo được đánh giá bằng các thước 
	đo xác suất và xếp hạng như Brier score, log-loss, calibration và khả năng 
	thu hồi chữ số thực tế trong nhóm dự báo có xác suất cao.
	
	\item \textbf{Đánh giá chiến lược:} xác suất dự báo được chuyển thành các tập 
	số ứng viên và chiến lược lựa chọn vé. Hiệu quả được đánh giá thông qua số lần 
	trúng thưởng, tổng chi phí, tổng tiền thưởng, lợi nhuận tích lũy và tỷ suất 
	hoàn vốn (ROI), đồng thời được so sánh với chiến lược ngẫu nhiên hoặc các 
	đường cơ sở thích hợp.
\end{itemize}
